\section{Reuse the power loop with multiword operations}
The exponentiation rule is unchanged when a residue occupies four words. What changes
is where state lives across a call. The nested modular multiplication can overwrite
argument and temporary registers. Retain the modulus pointer, output pointer, base,
exponent and result in saved registers or assigned stack slots. Save \texttt{ra} before
the first \texttt{jal}: a nested call replaces its old return address.

One possible private frame occupies 96 bytes. It contains sixteen-byte objects for
current base, remaining exponent, result and temporary output, followed by full-width
saved registers. This is a local implementation choice, not a requirement on equivalent
solutions. Copy caller objects word by word before using aligned doubleword accesses
inside the frame. Initialize the result as \(1\bmod N\), so exponent zero and modulus
one have their specified behavior.

For the right-to-left loop, inspect the low bit of the private exponent. When it is
one, call modular multiplication with the private result and base, writing into a
separate temporary object. Only after the call returns successfully should that object
replace the private result. In this layout result begins at stack offset 32, base at
offset zero, temporary output at offset 48, and \texttt{s0} retains the modulus pointer.

\Needspace{16\baselineskip}
\MARSCodeLabel{Multiply into separate private storage before replacing the accumulated result.}
\begin{minted}[samepage=true]{mars-idea}
daddiu $a0, $sp, 32
    move $a1, $sp
    move $a2, $s0
    daddiu $a3, $sp, 48
    jal rsa_mulmod128
    nop
    ld $t0, 48($sp)
    ld $t1, 56($sp)
    sd $t0, 32($sp)
    sd $t1, 40($sp)
\end{minted}

Shift the entire exponent right, using the transfer rule from modular multiplication.
If any exponent bits remain, square the current base modulo \(N\), again using a
separate temporary. Do not accidentally square the result in this right-to-left
variant: the result and current base have different roles.

A left-to-right variant is equally valid. Begin with \(R=1\bmod N\), visit exponent
bits from high to low, square \(R\), and multiply it by the original reduced base when
the next bit is one. After a prefix \(p\), the invariant is \(R\equiv\text{base}^{p}\).
Appending bit \(b\) changes the prefix to \(2p+b\), so the square followed by an optional
multiply produces the next value. For binary exponent \(10001\), the results after
its five significant bits are 65, 992, 1232, 1547 and 2790 modulo 3233.

These scan directions have different intermediate states. A grader for the completed
function should compare its interface result and state obligations rather than require
the right-to-left register trace. Guided observations can name a particular variant
when teaching that variant's invariant.
